\documentclass{article}%
\usepackage[T1]{fontenc}%
\usepackage[utf8]{inputenc}%
\usepackage{lmodern}%
\usepackage{textcomp}%
\usepackage{lastpage}%
\usepackage{geometry}%
\geometry{margin=2.2cm}%
%
\usepackage{amsmath}%
\usepackage{amssymb}%
\usepackage{booktabs}%
\usepackage{longtable}%
\usepackage{array}%
\usepackage{hyperref}%
\hypersetup{hidelinks}%
\title{IME 2027 {-} Química (claude\_haiku\_5\_5) x Poliedro\\[4pt]\large Comparação com o gabarito Poliedro Resolve\thanks{Gerado com gpt-resolve: \url{https://github.com/lgabs/gpt-resolve}}}%
\date{10/11/2026}%
%
\begin{document}%
\normalsize%
\maketitle%
\section*{Resumo}\textbf{Nota média: 9,5/10} (22 de 23 itens com a resolução correta; cada item vale o mesmo dentro da questão).\par\medskip%
\begin{longtable}{@{}c c p{0.78\textwidth}@{}}\toprule Questão & Nota & Resumo \\ \midrule
1 & 10,0 & A resolução e o gabarito do Poliedro concordam. Ambos trazem as mesmas reações (MnO2 + 4 HCl $\to$ MnCl2 + Cl2 + 2 H2O e 2 NaBr + Cl2 $\to$ 2 NaCl + Br2), a proporção 1 MnO2 : 2 NaBr e a mesma massa: 0,870 g na resolução e 0,87 g no gabarito. Refiz a conta com a Tabela da prova e obtive 0,8703 g, então os dois estão certos. Há apenas observações menores de redação e hipóteses implícitas razoáveis (MnO2 puro e HCl suficiente). \\
2 & 5,0 & O item b está certo e equivale ao gabarito. A resolução mostra que W2 + W4 = 0 pela 1ª lei com Q = 0 e porque U é função de estado, e chega a W\_tot = W1 + W3. As fórmulas diferem do gabarito só por um sinal global, porque a resolução usa o trabalho feito pelo gás e o gabarito, o trabalho feito sobre o gás. No item a, a sequência das etapas e a ordem das isotermas (T3 > T1) estão certas. Porém a resolução trata II e IV como irreversíveis e as desenha como cordas tracejadas, sem as adiabáticas do ciclo de Carnot que o gabarito e o enunciado esperam. Os estados escolhidos também não são compatíveis com adiabáticas reversíveis. Esse esboço é coerente com a hipótese que a resolução declara, mas foge da resposta esperada e deve perder parte da nota do item a. \\
3 & 10,0 & A resolução concorda com o gabarito Poliedro em todos os itens. Ordens: 2 em SO2 e 0 em O3. Lei: v = k{[}SO2{]}$^2$. Constante: k = 1,888 L$\cdot$mol$^{-}$$^{1}$$\cdot$s$^{-}$$^{1}$. Energia de ativação: Ea $\approx$ 4,93 kJ/mol, que a resolução arredonda para 4,9$\times$10$^3$ J/mol na resposta final. Refiz as contas e confirmei os valores: Ea/R $\approx$ 594 K e Ea $\approx$ 4930--4935 J/mol, conforme se use ln 4 ou ln 4,004. A única diferença no item d) é de arredondamento. A transcrição do enunciado confere com o original. Não há erro nem na resolução nem no gabarito. \\
4 & 10,0 & A resolução está correta e bate com o gabarito Poliedro nos itens a, b e d. No item d, refiz a conta e obtive i = 0,58, W = 0,1156 mol/kg, $\Delta$T = 0,343 K e Tf = 5,11 $^\circ$C. No item c, a forma é diferente, mas as respostas são equivalentes: o gabarito dá Kc = 16,4$\cdot$V, porque o enunciado não informa o volume, e a resolução declara a hipótese V = 1 L e obtém 16,4 L/mol, que é o mesmo valor. O único ponto para levar ao professor é se a banca aceita o valor numérico de Kc com essa hipótese. \\
5 & 10,0 & A resolução e o gabarito do Poliedro concordam: a substância é o n{-}heptano. Os dois chegam a $\Delta$n\_gás = 450$\cdot$16,6/(8,3$\cdot$300) = +3, que corresponde a y = 16 em CxHy, já que $\Delta$n = y/4 - 1. Conferi todas as contas e equações de combustão em Python, e a convenção de sinais está correta. A resolução é mais completa que o gabarito: testa as quatro alternativas, faz a análise de sinais e de ordem de grandeza e declara a premissa de água gasosa. Não encontrei erros. Restam só observações de baixa importância sobre a idealização do enunciado: temperatura fixa em 27 $^\circ$C e valores absolutos hipotéticos. \\
6 & 10,0 & A resolução está correta e coincide com o gabarito Poliedro nos três itens. Conferi a transcrição contra a página 3 da prova e refiz as contas em Python. a) Esterificação 1:1:1:1, com a mesma equação do gabarito. b) $\alpha$(ácido acético) = 1,6/4,0 = 40\%. c) Kc = (1,6$\cdot$1,6)/(2,4$\cdot$0,4) = 8/3 $\approx$ 2,67; o volume não entra porque $\Delta$n = 0. Não há divergência a resolver; só ficam observações menores sobre a redação do item b e sobre hipóteses implícitas. \\
7 & 10,0 & A resolução está correta e concorda com o gabarito Poliedro em todos os itens: $\alpha$ = 1, $\beta$ = 2, $\gamma$ = -1; k = 1,25$\times$10$^{-}$$^2$ L$\cdot$mol$^{-}$$^{1}$$\cdot$s$^{-}$$^{1}$; v = k{[}A{]}{[}B{]}$^2$/{[}C{]}; ordem global 2. Os dois supõem o mesmo erro de digitação no $\nu$$_{0}$ do experimento 4 (1,28$\times$10$^{-}$$^{6}$ $\to$ 1,28$\times$10$^{-}$$^3$). Refazendo as contas, a hipótese se confirma: o valor corrigido dá solução única (1, 2, -1), e o valor impresso leva a expoentes irracionais e a um k 10$^3$ vezes menor. O ponto a levar ao professor é a possível errata ou anulação do dado do experimento 4. \\
8 & 10,0 & A resolução está correta e bate com o gabarito do Poliedro nos dois itens. a) Y é o 3,4{-}dimetil{-}hexano meso. Confirmei com o produto misto que o desenho em zigue{-}zague dá (3R,4S) e é de fato o meso. b) X é o 2{-}clorobutano, e a observação de que ele precisa ser racêmico é correta. Não achei erro na resolução nem no gabarito; só há observações de redação e uma ambiguidade leve no enunciado sobre Y ser 'composto único'. \\
9 & 10,0 & A resolução coincide com o gabarito Poliedro e está correta. Pela crioscopia, M = 5,12$\times$53/(1,28$\times$2) = 106 g/mol, conta conferida. Entre os aromáticos com M = 106, o Tollens negativo exclui o benzaldeído. Só o p{-}xileno (1,4{-}dimetilbenzeno) dá um único produto de monohalogenação no anel. Há apenas imprecisões menores na enumeração de candidatos, que não alteram a resposta: a afirmação sobre anéis fundidos esquece o benzociclobuteno (104 g/mol), e o critério da tabela muda no caso da N{-}metilanilina. Nada que afete a nota. \\
10 & 10,0 & A resolução está certa e é equivalente ao gabarito Poliedro. As duas usam N\_Pb/N\_Th = e\^{}\{$\lambda$t\} - 1 = 0,024 com a meia{-}vida do $^2$$^3$$^2$Th. O gabarito (t/t$\frac12$ = 0,03463, t $\approx$ 4,88$\times$10$^{8}$ anos) usa a linearização ln(1,024) $\approx$ 0,024. A resolução apresenta esse mesmo 4,88$\times$10$^{8}$ e também o valor exato, $\approx$ 4,83$\times$10$^{8}$ anos, e declara a faixa 4,8--4,9$\times$10$^{8}$ anos. A diferença de \textasciitilde{}1\% está dentro da precisão de R = 0,024. Não há erro conceitual nem de conta em nenhum dos dois. \\
\bottomrule\end{longtable}%
\section*{Principais pontos de dúvida}%
\begin{itemize}%
\item%
\textbf{Q2} \textit{(solution)} Item a: a resolução se recusa a desenhar II e IV como adiabáticas (PV\^{}$\gamma$) e as representa como cordas tracejadas, porque as trata como irreversíveis. O enunciado ('muito rápida', sem troca de calor) e o item b apontam para o ciclo de Carnot com adiabáticas, que é o que o gabarito traz. Os estados da figura também não são compatíveis com adiabáticas reversíveis (V2V4 $\neq$ V1V3), então o desenho não tem a forma de Carnot. Um professor deve decidir quanto dessa interpretação seria aceito. É provável que custe parte da nota do item a.%
\item%
\textbf{Q4} \textit{(statement)} Item c: o enunciado não dá volume nem concentração, então Kc = {[}D{]}/{[}M{]}$^2$ depende do volume: Kc = 16,4$\cdot$V, ou 16,4/c0 em L/mol. O gabarito Poliedro deixa a resposta em função de V. A resolução adota 1 mol em 1 L e dá 16,4 L/mol, com a hipótese declarada. Vale perguntar ao professor se a correção oficial aceita o valor numérico sob essa hipótese ou se exige que a dependência do volume seja discutida.%
\item%
\textbf{Q7} \textit{(statement)} No experimento 4, o $\nu$$_{0}$ impresso (1,28$\times$10$^{-}$$^{6}$) não é compatível com nenhuma lei de potência de expoentes racionais. Tomado ao pé da letra, dá $\alpha$ $\approx$ -1,49, $\beta$ $\approx$ 1,00, $\gamma$ $\approx$ 0,49. Já 1,28$\times$10$^{-}$$^3$ é exatamente o valor previsto por (1, 2, -1): 1,25$\times$10$^{-}$$^{6}$ $\times$ 1024. A resolução e o gabarito supõem o mesmo erro de digitação. Vale conferir se o IME publicou errata ou anulou o item. Para os itens a) e d), é preciso definir como corrigir quem tomou o dado ao pé da letra.%
\end{itemize}%
\section*{Questão 1}%
\href{https://poliedroresolve.sistemapoliedro.com.br/vestibulares/ime/2027/ime-2-fase-dia-3-28-10-2026-2-fase-dia-3/dissertativa-1-quimica-2-fase-dia-3-ime-2027}{Resolução do gabarito}\par\smallskip%
\begin{longtable}{@{}p{0.06\textwidth} p{0.33\textwidth} p{0.33\textwidth} p{0.1\textwidth} p{0.08\textwidth}@{}}\toprule Item & Resolução & Gabarito & Veredito & Certo \\ \midrule
unico & MnO2 + 4 HCl $\to$ MnCl2 + Cl2 + 2 H2O; 2 NaBr + Cl2 $\to$ 2 NaCl + Br2 (global: MnO2 + 4 HCl + 2 NaBr $\to$ MnCl2 + 2 NaCl + Br2 + 2 H2O); m(MnO2) $\approx$ 0,870 g & MnO2 + 4HCl $\to$ MnCl2 + Cl2 + 2H2O; Cl2 + 2NaBr $\to$ 2NaCl + Br2; m(MnO2) = 0,87 g & igual & ambos \\
\bottomrule\end{longtable}%
\begin{itemize}%
\item%
\textbf{unico)} Refiz a conta com as massas molares da Tabela da prova (Na 22,990; Br 79,904; Mn 54,938; O 15,999): n(NaBr) = 2,06/102,894 = 2,002$\times$10$^{-}$$^2$ mol; n(MnO2) = n(Cl2) = n(NaBr)/2 = 1,001$\times$10$^{-}$$^2$ mol; m = 1,001$\times$10$^{-}$$^2$ $\times$ 86,936 = 0,8703 g. Com massas inteiras (103 e 87) a conta dá exatamente 0,87 g. As duas reações são iguais às do gabarito e estão balanceadas. A diferença entre 0,870 e 0,87 é só de algarismos significativos. Conferi também as verificações da resolução: a massa de Br é 1,600 g, e com Br = 80,00 a massa vai a 0,8694 g, diferença menor que 0,001 g.%
\end{itemize}%
\paragraph{Pontos de dúvida}%
\begin{itemize}%
\item%
\textit{(statement)} O enunciado fala em pirolusita, mas não informa o teor de MnO2 no minério nem a quantidade de HCl. A resolução e o gabarito calculam a massa de MnO2 puro com HCl suficiente, que é a leitura natural da pergunta ('massa de óxido de manganês IV'). A resolução declara essas hipóteses explicitamente.%
\item%
\textit{(solution)} Na parte (b) a resolução escreve '(cada Br2 perde 2 e$^{-}$)'. O mais preciso seria 'os 2 Br$^{-}$ que formam cada Br2 perdem 2 e$^{-}$'. Fora isso, a hipótese (iii), sobre a formação de BrCl com excesso de Cl2, é um comentário a mais que não interfere no resultado. São só questões de redação.%
\end{itemize}%
\section*{Questão 2}%
\href{https://poliedroresolve.sistemapoliedro.com.br/vestibulares/ime/2027/ime-2-fase-dia-3-28-10-2026-2-fase-dia-3/dissertativa-2-quimica-2-fase-dia-3-ime-2027}{Resolução do gabarito}\par\smallskip%
\begin{longtable}{@{}p{0.06\textwidth} p{0.33\textwidth} p{0.33\textwidth} p{0.1\textwidth} p{0.08\textwidth}@{}}\toprule Item & Resolução & Gabarito & Veredito & Certo \\ \midrule
a & Esboço P$\times$V com as isotermas quase{-}estáticas 1$\to$2 (T1) e 3$\to$4 (T3 > T1) desenhadas como curvas. As etapas II e IV aparecem só como segmentos retos tracejados ligando os extremos, porque a resolução as trata como irreversíveis e diz que não têm curva P$\times$V. Ela recusa explicitamente desenhar a adiabática PV\^{}$\gamma$. Ordenação: V3 < V2 < V1, V3 < V4 < V1, P1 mínima e P3 máxima. Os valores da figura (V1=12, V2=5, V3=3, V4=4; RT1=60, RT3=120) não satisfazem V2V4 = V1V3. & Gráfico P$\times$V do ciclo 1$\to$2 (I, isoterma T1), 2$\to$3 (II, adiabática), 3$\to$4 (III, isoterma T3), 4$\to$1 (IV, adiabática). É o ciclo de Carnot, com as etapas II e IV desenhadas como adiabáticas, mais íngremes que as isotermas. & diferente & gabarito \\
b & W2 = -$\Delta$U2 = -{[}U(T3) - U(T1){]} e W4 = -$\Delta$U4 = -{[}U(T1) - U(T3){]}. Logo W2 + W4 = 0, também justificado por U ser função de estado com $\Delta$U1 = $\Delta$U3 = 0. Portanto W\_tot = W1 + W3 = RT1 ln(V2/V1) + RT3 ln(V4/V3), na convenção W = trabalho realizado pelo gás. & W2 = nCv(T3 - T1) e W4 = nCv(T1 - T3) = -W2, logo W\_tot = W1 + W3, com W1 = -nRT1 ln(V2/V1) e W3 = -nRT3 ln(V4/V3). Convenção: W = trabalho realizado sobre o gás. & equivalente & ambos \\
\bottomrule\end{longtable}%
\begin{itemize}%
\item%
\textbf{a)} O enunciado confere com o original (pág. 2). A sequência das etapas, o sentido de cada uma e o fato de a isoterma III ficar acima da isoterma I (T3 > T1) estão certos nas duas respostas. A divergência está em como II e IV são representadas. O enunciado diz que essas etapas são 'muito rápidas' só para justificar Q = 0. Na leitura usual de livro{-}texto, e claramente a pretendida, já que o item b é o passo clássico da dedução do ciclo de Carnot, elas são adiabáticas e se desenham como curvas PV\^{}$\gamma$ = cte, mais íngremes que as isotermas. A resolução preferiu uma leitura irreversível estrita: trocou as adiabáticas por cordas tracejadas e escolheu estados incompatíveis com adiabáticas reversíveis (V2V4 = 20 $\neq$ V1V3 = 36). Com isso o esboço não tem a forma de um ciclo de Carnot, e a corda IV chega a passar acima da isoterma T3 logo depois do estado 4. Sob a hipótese que a própria resolução declara, o esboço é coerente. Conferi a 2ª lei: S3 $\geq$ S2 e S1 $\geq$ S4 exigem 0,737R $\leq$ Cv $\leq$ 1,585R, o que vale para gás monoatômico. Mesmo assim, ele foge da resposta esperada (adiabáticas desenhadas) e deve perder parte da nota. Por isso considero o gabarito certo neste item.%
\item%
\textbf{b)} A conclusão é a mesma: W2 + W4 = 0, logo W\_tot = W1 + W3. As expressões diferem por um sinal global porque as convenções diferem. O gabarito usa $\Delta$U = Q + W, com W feito sobre o gás. Daí W2 = $\Delta$U2 = nCv(T3 - T1) > 0 e W1 = -nRT1 ln(V2/V1) > 0 na compressão. A resolução usa W feito pelo gás: W2 = -$\Delta$U2 < 0 e W1 = RT1 ln(V2/V1) < 0. Cada parcela de uma é exatamente o oposto da correspondente na outra, e as duas são internamente coerentes. O argumento da resolução é até mais geral: não exige Cv constante, só U = U(T) e o fato de U ser função de estado. Conferi também a verificação dela para o caso reversível: somando as adiabáticas, V1V3 = V2V4 e W\_tot = R(T3 - T1) ln(V1/V2). Está correta.%
\end{itemize}%
\paragraph{Pontos de dúvida}%
\begin{itemize}%
\item%
\textit{(solution)} Item a: a resolução se recusa a desenhar II e IV como adiabáticas (PV\^{}$\gamma$) e as representa como cordas tracejadas, porque as trata como irreversíveis. O enunciado ('muito rápida', sem troca de calor) e o item b apontam para o ciclo de Carnot com adiabáticas, que é o que o gabarito traz. Os estados da figura também não são compatíveis com adiabáticas reversíveis (V2V4 $\neq$ V1V3), então o desenho não tem a forma de Carnot. Um professor deve decidir quanto dessa interpretação seria aceito. É provável que custe parte da nota do item a.%
\item%
\textit{(statement)} Ambiguidade do enunciado: 'compressão muito rápida' pode ser lida como adiabática quase{-}estática (rápida em relação à troca de calor, mas lenta em relação ao equilíbrio de pressão), que é a leitura usual, ou como processo irreversível. A conclusão do item b (W2 + W4 = 0) vale nas duas leituras. Só o gráfico do item a depende da escolha.%
\item%
\textit{(key)} O gabarito usa a convenção 'trabalho realizado sobre o gás' ($\Delta$U = Q + W), comum em Química, e daí W1 = -nRT1 ln(V2/V1). A resolução usa 'trabalho realizado pelo gás'. As expressões diferem por um sinal global, mas são equivalentes. O gabarito também supõe Cv constante (W2 = nCv$\Delta$T), o que não é necessário para a demonstração.%
\item%
\textit{(solution)} A resolução diz que a figura 'mostra um caso possível', mas não declara o Cv. Com T3/T1 = 2 e os volumes usados, a 2ª lei nas etapas irreversíveis (S3 $\geq$ S2 e S1 $\geq$ S4) exige 0,737R $\leq$ Cv $\leq$ 1,585R. A figura vale para gás monoatômico, mas violaria a 2ª lei na etapa IV para gás diatômico (Cv = 2,5R).%
\end{itemize}%
\section*{Questão 3}%
\href{https://poliedroresolve.sistemapoliedro.com.br/vestibulares/ime/2027/ime-2-fase-dia-3-28-10-2026-2-fase-dia-3/dissertativa-3-quimica-2-fase-dia-3-ime-2027}{Resolução do gabarito}\par\smallskip%
\begin{longtable}{@{}p{0.06\textwidth} p{0.33\textwidth} p{0.33\textwidth} p{0.1\textwidth} p{0.08\textwidth}@{}}\toprule Item & Resolução & Gabarito & Veredito & Certo \\ \midrule
a & Ordem 2 em relação ao SO2 e ordem 0 em relação ao O3 (ordem global 2) & Ordem 2 em SO2, ordem 0 em O3 & igual & ambos \\
b & v = k{[}SO2{]}\^{}2 & v = k{[}SO2{]}\^{}2 & igual & ambos \\
c & k = 1,888 $\approx$ 1,89 L$\cdot$mol$^{-}$$^{1}$$\cdot$s$^{-}$$^{1}$ (300 K) & k = 1,888 L$\cdot$mol$^{-}$$^{1}$$\cdot$s$^{-}$$^{1}$ & igual & ambos \\
d & Ea $\approx$ 4,9$\times$10$^3$ J$\cdot$mol$^{-}$$^{1}$ (cálculo: 4,93$\times$10$^3$ J/mol; $\approx$1,18$\times$10$^3$ cal/mol), energia de ativação aparente & Ea = 594R $\approx$ 4,93 kJ/mol & equivalente & ambos \\
\bottomrule\end{longtable}%
\begin{itemize}%
\item%
\textbf{a)} Exp. 1 e 2: v não muda quando {[}O3{]} cai à metade, logo b = 0. Exp. 2 e 3: {[}SO2{]} triplica e v fica 9 vezes maior (1,062/0,118 = 9), logo a = 2. A resolução ainda confere o par 1--3, que dá o mesmo resultado.%
\item%
\textbf{b)} É consequência direta do item a). A unidade L$\cdot$mol$^{-}$$^{1}$$\cdot$s$^{-}$$^{1}$ indicada pela resolução está correta para uma lei de 2ª ordem.%
\item%
\textbf{c)} Pelos exp. 1 e 2: 0,118/0,0625 = 1,888. Pelo exp. 3: 1,062/0,5625 = 1,888. A resolução dá o valor exato e o arredondado. A diferença é só de arredondamento.%
\item%
\textbf{d)} Refiz a conta: ln(7,56/1,888) = ln(4,0042) = 1,3874 e 1/300 - 1/1000 = 7/3000 K$^{-}$$^{1}$. Isso dá Ea/R = 594,6 K e Ea = 8,3 $\times$ 594,6 $\approx$ 4935 J/mol. O gabarito usa ln 4 = 2 $\times$ 0,693 = 1,386, o que dá Ea/R = 594 exatamente e Ea = 594 $\times$ 8,3 = 4930 J/mol $\approx$ 4,93 kJ/mol. A resolução chega a 4,93$\times$10$^3$ J/mol no desenvolvimento e arredonda a resposta final para 4,9$\times$10$^3$ (2 algarismos significativos, por causa de R = 8,3). É o mesmo resultado, só com arredondamento diferente. O valor em cal (1,18$\times$10$^3$ cal/mol) também confere: 1,987 $\times$ 594,6 $\approx$ 1181.%
\end{itemize}%
\paragraph{Pontos de dúvida}%
\begin{itemize}%
\item%
\textit{(solution)} No item d), a resposta final da resolução foi arredondada para 4,9$\times$10$^3$ J/mol (2 algarismos significativos), e o gabarito dá 4,93 kJ/mol. O valor 4,93$\times$10$^3$ aparece explicitamente no desenvolvimento, então isso não deve custar pontos. A pequena diferença em Ea/R (594,6 contra 594) vem de a resolução usar ln(4,004) e o gabarito arredondar para ln 4 = 2 ln 2.%
\item%
\textit{(statement)} O enunciado dá k = 7,56 a 1000 K apenas 'em módulo', sem unidade. A resolução e o gabarito supõem implicitamente a mesma unidade de k a 300 K (L$\cdot$mol$^{-}$$^{1}$$\cdot$s$^{-}$$^{1}$), ou seja, mesma lei de velocidade e mesmo mecanismo a 1000 K. Também supõem que a luz UV continue igual. Sem isso, a razão k2/k1 não teria sentido.%
\item%
\textit{(solution)} Aplicar Arrhenius a uma reação fotoquímica com apenas dois pontos de temperatura dá uma Ea aparente. A resolução declara essa limitação. Os comentários sobre mecanismo, saturação e rendimento quântico são especulativos, mas não afetam as respostas pedidas.%
\item%
\textit{(solution)} A ordem zero em O3 só foi verificada na faixa de 0,20 a 0,40 mol/L, ressalva que a própria resolução faz. É uma hipótese razoável e não muda a resposta.%
\end{itemize}%
\section*{Questão 4}%
\href{https://poliedroresolve.sistemapoliedro.com.br/vestibulares/ime/2027/ime-2-fase-dia-3-28-10-2026-2-fase-dia-3/dissertativa-4-quimica-2-fase-dia-3-ime-2027}{Resolução do gabarito}\par\smallskip%
\begin{longtable}{@{}p{0.06\textwidth} p{0.33\textwidth} p{0.33\textwidth} p{0.1\textwidth} p{0.08\textwidth}@{}}\toprule Item & Resolução & Gabarito & Veredito & Certo \\ \midrule
a & Ligações de hidrogênio, duas pontes O--H$\cdot$$\cdot$$\cdot$O=C por dímero, que fecham um anel de 8 membros. & Duas ligações de hidrogênio entre as carboxilas. & igual & ambos \\
b & Dímero cíclico C4H8O4 com duas pontes de hidrogênio tracejadas (O--H$\cdot$$\cdot$$\cdot$O=C) e anel de 8 membros. A figura traz os pares não ligantes nos O, e a contagem de 48 elétrons de valência está conferida. & Dímero cíclico com duas ligações de H (estrutura em imagem). & igual & ambos \\
c & Kc = {[}D{]}/{[}M{]}$^2$ = 0,42/0,16$^2$ $\approx$ 16,4 L$\cdot$mol$^{-}$$^{1}$, assumindo de forma explícita 1 mol de monômero inicial em 1 L (c0 = 1 mol/L). & Kc = 0,42V/0,16$^2$ = 16,4$\cdot$V (V em L). Segundo o gabarito, o enunciado não dá o volume, então não há valor numérico. & equivalente & ambos \\
d & Mesma razão D/M = 2,625: {[}M{]} = 0,016 e {[}D{]} = 0,042 mol/L; i = 0,58; W $\approx$ 0,1156 mol/kg; $\Delta$Tf $\approx$ 0,343 K; Tf $\approx$ 5,11 $^\circ$C. & i = 0,58; W = 0,1156 mol/kg; $\Delta$Tc = 0,343 $^\circ$C; T = 5,11 $^\circ$C. & igual & ambos \\
\bottomrule\end{longtable}%
\begin{itemize}%
\item%
\textbf{a)} A resolução e o gabarito dão a mesma resposta: duas ligações de hidrogênio entre a hidroxila de uma molécula e a carbonila da outra. A resolução também cita, corretamente, as forças dipolo{-}dipolo e de dispersão como interações secundárias.%
\item%
\textbf{b)} A estrutura está correta. O anel C1--O(H)--H$\cdot$$\cdot$$\cdot$O=C2--O(H)--H$\cdot$$\cdot$$\cdot$O=C1 tem 8 membros, cada C forma 4 ligações, cada O tem 2 pares não ligantes e a contagem dá 24 elétrons por monômero e 48 no total. Os detalhes de desenho (CH3 condensado, pares da carbonila não voltados para a ponte) são só estéticos.%
\item%
\textbf{c)} Refiz a conta: com 1 mol inicial em volume V, {[}D{]} = 0,42/V e {[}M{]} = 0,16/V, logo Kc = (0,42/V)/(0,16/V)$^2$ = 16,406$\cdot$V. A resolução escreveu que o volume não foi dado e adotou V = 1 L. Com V = 1, a expressão do gabarito vira exatamente o valor da resolução, por isso as respostas são equivalentes. O gabarito é mais rigoroso, porque deixa a dependência de V explícita. A resolução dá um caso particular, com a hipótese declarada, que provavelmente é o valor que a banca esperava.%
\item%
\textbf{d)} Refiz a conta. M + 2D = 0,10 com D = 2,625M dá M = 0,016 e D = 0,042, logo i = 0,058/0,10 = 0,58. Massa de ácido 6,0 g ocupa 5,7 mL; benzeno 0,9943 L $\times$ 0,87 = 0,8650 kg; W = 0,11560 mol/kg. $\Delta$T = 0,58 $\times$ 5,12 $\times$ 0,1156 = 0,3433 K, então Tf = 5,107 $\approx$ 5,11 $^\circ$C. Usar M = 60 ou 60,05 g/mol não muda o resultado. Desprezando o volume do soluto, W = 0,1149 e Tf = 5,109 $\approx$ 5,11 $^\circ$C. Também conferi as verificações da resolução: com i = 1, $\Delta$T = 0,592 K; com Kc fixo em 0,10 mol/L, a razão D/M seria 0,69.%
\end{itemize}%
\paragraph{Pontos de dúvida}%
\begin{itemize}%
\item%
\textit{(statement)} Item c: o enunciado não dá volume nem concentração, então Kc = {[}D{]}/{[}M{]}$^2$ depende do volume: Kc = 16,4$\cdot$V, ou 16,4/c0 em L/mol. O gabarito Poliedro deixa a resposta em função de V. A resolução adota 1 mol em 1 L e dá 16,4 L/mol, com a hipótese declarada. Vale perguntar ao professor se a correção oficial aceita o valor numérico sob essa hipótese ou se exige que a dependência do volume seja discutida.%
\item%
\textit{(statement)} A frase 'cada mol de monômero gera 0,42 mol de dímero' admite outra leitura: razão D/M = 0,42 no equilíbrio. Por ela, o item d teria i $\approx$ 0,77, e não 0,58. A resolução e o gabarito usam a mesma leitura (1 mol inicial $\to$ 0,42 mol de D e 0,16 mol de M), que é a mais natural.%
\item%
\textit{(statement)} Item d: ao exigir a 'mesma relação dímero/monômero' em 0,10 mol/L, o enunciado contradiz um Kc constante, que daria D/M $\approx$ 0,69 nessa concentração. A resolução aponta isso e segue o enunciado, assim como o gabarito.%
\item%
\textit{(solution)} Item d: para chegar à molalidade, a resolução supõe volumes aditivos e desconta o volume do ácido (densidade 1,05). É a mesma conta do gabarito (W = 0,1156), e desprezar esse volume dá o mesmo Tf $\approx$ 5,11 $^\circ$C.%
\item%
\textit{(solution)} Item b: a figura escreve o CH3 de forma condensada, sem mostrar as ligações C--H, e os pares não ligantes do O da carbonila não ficam voltados para a ponte de hidrogênio. Isso não afeta a correção da estrutura, que tem 48 elétrons, octetos completos e anel de 8 membros.%
\end{itemize}%
\section*{Questão 5}%
\href{https://poliedroresolve.sistemapoliedro.com.br/vestibulares/ime/2027/ime-2-fase-dia-3-28-10-2026-2-fase-dia-3/dissertativa-5-quimica-2-fase-dia-3-ime-2027}{Resolução do gabarito}\par\smallskip%
\begin{longtable}{@{}p{0.06\textwidth} p{0.33\textwidth} p{0.33\textwidth} p{0.1\textwidth} p{0.08\textwidth}@{}}\toprule Item & Resolução & Gabarito & Veredito & Certo \\ \midrule
unico & n{-}heptano (C7H16). $\Delta$n = y/4 - 1; $\Delta$H - $\Delta$U = C($\Delta$T\_V - $\Delta$T\_P) = 450$\cdot$16,6 = 7470 J/mol = $\Delta$n$\cdot$RT, logo $\Delta$n = 7470/2490 = +3, y = 16. C7H16 + 11 O2 $\to$ 7 CO2 + 8 H2O. & $\Delta$n\_gás = 450$\cdot$16,6/(8,3$\cdot$300) = 3, logo n{-}heptano (C7H16 + 11 O2 $\to$ 7 CO2 + 8 H2O). & igual & ambos \\
\bottomrule\end{longtable}%
\begin{itemize}%
\item%
\textbf{unico)} Mesma conta e mesma conclusão. Refiz tudo: 450$\cdot$16,6 = 7470 J e 8,3$\cdot$300 = 2490 J/mol, o que dá $\Delta$n = 3,000. As equações da tabela também conferem. Acetileno: 3,5 $\to$ 3, $\Delta$n = -0,5, diferença de -2,77 K. p{-}xileno: 11,5 $\to$ 13, $\Delta$n = +1,5, diferença de 8,30 K. n{-}heptano: 12 $\to$ 15, $\Delta$n = +3, diferença de 16,6 K. Propeno: 5,5 $\to$ 6, $\Delta$n = +0,5, diferença de 2,77 K. A convenção de sinais está certa: $\Delta$U = -C$\Delta$T\_V e $\Delta$H = -C$\Delta$T\_P, então $\Delta$H - $\Delta$U = C($\Delta$T\_V - $\Delta$T\_P) > 0, e $\Delta$n > 0 é coerente com (v). O enunciado transcrito confere com a página 3 do PDF.%
\end{itemize}%
\paragraph{Pontos de dúvida}%
\begin{itemize}%
\item%
\textit{(statement)} O enunciado diz que a reação ocorreu 'à mesma temperatura de 27 $^\circ$C', mas os calorímetros se aquecem ($\Delta$T $\neq$ 0). A leitura usual, adotada pela resolução e pelo gabarito, é calcular $\Delta$nRT com T = 300 K e desprezar a correção de Kirchhoff. Os valores absolutos são hipotéticos: cerca de 4,8 MJ aqueceriam um calorímetro de 450 J/K em uns 10$^{4}$ K. A resolução comenta isso, e a identificação não muda, porque só depende da diferença $\Delta$T\_V - $\Delta$T\_P.%
\item%
\textit{(statement)} A premissa (ii), todos os produtos gasosos, inclusive a H2O, é essencial. Com água líquida, nenhuma alternativa daria $\Delta$n = +3. A resolução observa isso explicitamente. O gabarito usa a mesma premissa sem comentar.%
\item%
\textit{(solution)} Para que a conta dê o número inteiro exato, é preciso usar R = 8,3 J$\cdot$mol$^{-}$$^{1}$$\cdot$K$^{-}$$^{1}$, e não 8,314. A resolução justifica a escolha pela Folha de Dados. Com 8,314 o resultado seria $\Delta$n $\approx$ 2,995, e a conclusão não muda.%
\end{itemize}%
\section*{Questão 6}%
\href{https://poliedroresolve.sistemapoliedro.com.br/vestibulares/ime/2027/ime-2-fase-dia-3-28-10-2026-2-fase-dia-3/dissertativa-6-quimica-2-fase-dia-3-ime-2027}{Resolução do gabarito}\par\smallskip%
\begin{longtable}{@{}p{0.06\textwidth} p{0.33\textwidth} p{0.33\textwidth} p{0.1\textwidth} p{0.08\textwidth}@{}}\toprule Item & Resolução & Gabarito & Veredito & Certo \\ \midrule
a & CH3COOH(l) + C2H5OH(l) ? CH3COOC2H5(l) + H2O(l), coeficientes 1:1:1:1 (esterificação de Fischer) & CH3COOH(l) + C2H5OH(l) ? CH3COOC2H5(l) + H2O(l) & igual & ambos \\
b & $\alpha$\_ac = 1,6/4,0 = 0,40 = 40\% & $\alpha$ = 1,6/4 = 0,4 (40\%) & igual & ambos \\
c & Kc = (1,6$\cdot$1,6)/(2,4$\cdot$0,4) = 8/3 $\approx$ 2,7 (adimensional, pois $\Delta$n = 0) & Kc = (1,6$\cdot$1,6)/(2,4$\cdot$0,4) = 8/3 & igual & ambos \\
\bottomrule\end{longtable}%
\begin{itemize}%
\item%
\textbf{a)} É a mesma equação, com os mesmos estados físicos e coeficientes unitários. O balanço atômico confere: C 4=4, H 10=10, O 3=3.%
\item%
\textbf{b)} n0(ácido) = 240/60 = 4,0 mol e o consumo é de 1,6 mol, logo $\alpha$ = 0,40. A resolução também mostra que 80\% é a conversão do etanol (limitante) e não o valor pedido. Com as massas molares exatas da tabela, $\alpha$ $\approx$ 0,4003.%
\item%
\textbf{c)} Refiz a conta: 2,56/0,96 = 8/3 $\approx$ 2,67. Como $\Delta$n = 0, o volume se cancela. A conferência da resolução também está certa: 5x$^2$ - 48x + 64 = 0 tem raízes 1,6 e 8,0, e só 1,6 é aceitável. Com as massas molares exatas (60,05 e 46,07), Kc $\approx$ 2,69.%
\end{itemize}%
\paragraph{Pontos de dúvida}%
\begin{itemize}%
\item%
\textit{(statement)} A expressão 'grau de equilíbrio do sistema reacional para o ácido acético' é pouco usual. O gabarito e a resolução a leem como a fração de ácido que reagiu (40\%). Um aluno que respondesse 80\% (pelo etanol, que é o limitante) erraria o item, mas o 'para o ácido acético' do enunciado deixa a leitura de 40\% clara.%
\item%
\textit{(solution)} A resolução lê 'somente o equilíbrio líquido{-}líquido' como 'só fase líquida, sem equilíbrio líquido{-}vapor, mistura homogênea'. É uma leitura razoável e não muda nada nas contas. As hipóteses extras (catalisador H2SO4, nenhuma reação paralela, atividade aproximada pela concentração) são implícitas e razoáveis.%
\item%
\textit{(solution)} Massas molares arredondadas para 60 e 46 g/mol. Com os valores da tabela, Kc $\approx$ 2,69 em vez de 2,67; a diferença é desprezível e a resposta continua sendo 8/3 $\approx$ 2,7.%
\end{itemize}%
\section*{Questão 7}%
\href{https://poliedroresolve.sistemapoliedro.com.br/vestibulares/ime/2027/ime-2-fase-dia-3-28-10-2026-2-fase-dia-3/dissertativa-7-quimica-2-fase-dia-3-ime-2027}{Resolução do gabarito}\par\smallskip%
\begin{longtable}{@{}p{0.06\textwidth} p{0.33\textwidth} p{0.33\textwidth} p{0.1\textwidth} p{0.08\textwidth}@{}}\toprule Item & Resolução & Gabarito & Veredito & Certo \\ \midrule
a & $\alpha$ = 1, $\beta$ = 2, $\gamma$ = -1. Os expoentes saem dos experimentos 1 a 3 e do mecanismo ($\alpha$ = 1). O experimento 4 só serve de checagem, com a hipótese de que $\nu$$_{0}$ seja 1,28$\times$10$^{-}$$^3$. & Supondo erro de digitação no $\nu$$_{0}$ do exp. 4 (1,28$\cdot$10$^{-}$$^3$ em vez de 1,28$\cdot$10$^{-}$$^{6}$): $\alpha$ = 1, $\beta$ = 2, $\gamma$ = -1 & igual & ambos \\
b & k = 1,25$\times$10$^{-}$$^2$ L$\cdot$mol$^{-}$$^{1}$$\cdot$s$^{-}$$^{1}$ (= k$_{1}$k$_{2}$/k$_{3}$) & k = 1,25$\cdot$10$^{-}$$^2$ L$\cdot$mol$^{-}$$^{1}$$\cdot$s$^{-}$$^{1}$ & igual & ambos \\
c & $\nu$$_{0}$ = 1,25$\times$10$^{-}$$^2$$\cdot${[}A{]}{[}B{]}$^2$/{[}C{]} (na faixa dos dados). Também apresenta a forma mecanística $\nu$$_{0}$ = (k$_{1}$k$_{2}$/k$_{3}$){[}A{]}{[}B{]}$^2$/({[}C{]} + (k$_{2}$/k$_{3}$){[}B{]}). & v = k{[}A{]}{[}B{]}$^2${[}C{]}$^{-}$$^{1}$ & equivalente & ambos \\
d & Ordem global = 2 (1 em A, 2 em B, -1 em C) & ordem global = 2 & igual & ambos \\
\bottomrule\end{longtable}%
\begin{itemize}%
\item%
\textbf{a)} Refiz as contas. Com $\nu$$_{0}$,$_{4}$ = 1,28$\times$10$^{-}$$^3$, o sistema $\alpha$+2$\beta$+3$\gamma$=2, 2$\alpha$+$\beta$+4$\gamma$=0, 3$\alpha$+4$\beta$+$\gamma$=10 tem determinante 20 e solução única (1, 2, -1). Com o valor impresso (1,28$\times$10$^{-}$$^{6}$), a solução é (-1,491; 1,003; 0,495), com expoentes irracionais. Além disso, o k calculado pelo exp. 4 daria 1,25$\times$10$^{-}$$^{5}$, exatamente 10$^3$ vezes menor que nos exps. 1 a 3. Isso confirma o erro de potência de dez que a resolução e o gabarito supõem.%
\item%
\textbf{b)} k = $\nu$$_{0}${[}C{]}/({[}A{]}{[}B{]}$^2$) = 1,25$\times$10$^{-}$$^2$ nos exps. 1, 2 e 3, conferido numericamente. A unidade de uma lei de ordem global 2 é L$\cdot$mol$^{-}$$^{1}$$\cdot$s$^{-}$$^{1}$. A identificação k = k$_{1}$k$_{2}$/k$_{3}$ é um acréscimo coerente.%
\item%
\textbf{c)} A lei de potência da resolução é a mesma do gabarito, com o valor numérico de k já substituído. A expressão do estado estacionário é um complemento. Ela se reduz à lei do gabarito quando (k$_{2}$/k$_{3}$){[}B{]} ? {[}C{]}, condição que a resolução mostra valer nos dados. Ela não substitui a resposta pedida.%
\item%
\textbf{d)} 1 + 2 - 1 = 2.%
\end{itemize}%
\paragraph{Pontos de dúvida}%
\begin{itemize}%
\item%
\textit{(statement)} No experimento 4, o $\nu$$_{0}$ impresso (1,28$\times$10$^{-}$$^{6}$) não é compatível com nenhuma lei de potência de expoentes racionais. Tomado ao pé da letra, dá $\alpha$ $\approx$ -1,49, $\beta$ $\approx$ 1,00, $\gamma$ $\approx$ 0,49. Já 1,28$\times$10$^{-}$$^3$ é exatamente o valor previsto por (1, 2, -1): 1,25$\times$10$^{-}$$^{6}$ $\times$ 1024. A resolução e o gabarito supõem o mesmo erro de digitação. Vale conferir se o IME publicou errata ou anulou o item. Para os itens a) e d), é preciso definir como corrigir quem tomou o dado ao pé da letra.%
\item%
\textit{(solution)} Os experimentos 1 a 3 dão só duas equações para três incógnitas. A resolução fecha o sistema com o mecanismo, usando a aproximação do estado estacionário: $\alpha$ = 1 porque $\nu$ $\propto$ {[}A{]}. É uma hipótese adicional, declarada e razoável, mas não é obrigatória pelo enunciado. Sem o experimento 4 corrigido, os dados sozinhos não determinam os expoentes.%
\item%
\textit{(solution)} No item c), a resolução chama de 'lei completa' a expressão mecanística com denominador {[}C{]} + (k$_{2}$/k$_{3}$){[}B{]}, que não é uma lei de potência. O enunciado pede a lei na forma empírica $\nu$$_{0}$ = k{[}A{]}\^{}$\alpha${[}B{]}\^{}$\beta${[}C{]}\^{}$\gamma$, e a resolução também dá essa forma. Não há erro, só uma possível confusão de terminologia.%
\end{itemize}%
\section*{Questão 8}%
\href{https://poliedroresolve.sistemapoliedro.com.br/vestibulares/ime/2027/ime-2-fase-dia-3-28-10-2026-2-fase-dia-3/dissertativa-8-quimica-2-fase-dia-3-ime-2027}{Resolução do gabarito}\par\smallskip%
\begin{longtable}{@{}p{0.06\textwidth} p{0.33\textwidth} p{0.33\textwidth} p{0.1\textwidth} p{0.08\textwidth}@{}}\toprule Item & Resolução & Gabarito & Veredito & Certo \\ \midrule
a & Y = (3R,4S){-}3,4{-}dimetil{-}hexano, forma meso (C8H18), desenhado em zigue{-}zague: em C3, CH3 em cunha e H tracejado; em C4, CH3 tracejada e H em cunha & Y = forma meso do 3,4{-}dimetil{-}hexano (estrutura 3D em imagem) & igual & ambos \\
b & 2{-}clorobutano (cloreto de sec{-}butila), CH3CH2CHClCH3, usado como racemato (RS) & 2{-}clorobutano & igual & ambos \\
\bottomrule\end{longtable}%
\begin{itemize}%
\item%
\textbf{a)} Conferi o enunciado no PDF (p. 4, 8ª questão) e a transcrição está fiel. Na reação de Wurtz (2 RCl + 2 Na $\to$ R--R + 2 NaCl), Y = R--R tem metades iguais. Para Y ter 2 carbonos quirais, cada R precisa ter exatamente 1. A menor alquila saturada que cumpre isso é a sec{-}butila, e então Y é o 3,4{-}dimetil{-}hexano (C8). Entre seus estereoisômeros, só o meso (R,S) é uma substância única e opticamente inativa. Também conferi o desenho: calculei a quiralidade pelo produto misto, com as coordenadas da figura e o H pondo cunha/tracejado como $\pm$z. O resultado foi C3 = R e C4 = S. A inversão no ponto médio de C3--C4 leva a molécula nela mesma, então o desenho é de fato o meso. As prioridades CIP também estão certas: C4 > C2 > CH3 > H em C3, e o simétrico em C4.%
\item%
\textbf{b)} Mesmo nome IUPAC do gabarito. A resolução acrescenta que X tem de ser racêmico, já que o meso exige unir um fragmento R a um S. Isso está correto e não contradiz o gabarito. Ela também lembra que, entre os isômeros C4H9Cl, só o 2{-}clorobutano tem carbono quiral, o que também está correto.%
\end{itemize}%
\paragraph{Pontos de dúvida}%
\begin{itemize}%
\item%
\textit{(statement)} O enunciado chama Y de 'esse composto' único e inativo. Na prática, o 2{-}clorobutano racêmico em Wurtz daria uma mistura de meso com o par (R,R)/(S,S), que também é inativa no total. Pelo enunciado, Y é o meso. A resolução trata bem esse ponto, mas a ambiguidade está no texto da prova.%
\item%
\textit{(solution)} A resolução diz que o meso não desvia a luz 'pois as rotações de C3 e C4 se cancelam'. É uma explicação didática comum, mas imprecisa: o motivo rigoroso é a molécula ser aquiral (tem centro de inversão), e a própria resolução também demonstra isso.%
\item%
\textit{(solution)} Hipótese declarada: 'alquila' significa grupo saturado. Ela exclui alternativas como o 3{-}clorobut{-}1{-}eno, que daria um dieno C8 meso. É razoável e coerente com o termo 'cloreto de alquila'.%
\item%
\textit{(solution)} A resolução não comenta que, na prática, haletos secundários em Wurtz rendem pouco e têm reações paralelas (eliminação/desproporcionamento, dando butano e buteno). Isso não afeta a resposta pedida.%
\end{itemize}%
\section*{Questão 9}%
\href{https://poliedroresolve.sistemapoliedro.com.br/vestibulares/ime/2027/ime-2-fase-dia-3-28-10-2026-2-fase-dia-3/dissertativa-9-quimica-2-fase-dia-3-ime-2027}{Resolução do gabarito}\par\smallskip%
\begin{longtable}{@{}p{0.06\textwidth} p{0.33\textwidth} p{0.33\textwidth} p{0.1\textwidth} p{0.08\textwidth}@{}}\toprule Item & Resolução & Gabarito & Veredito & Certo \\ \midrule
unico & M = 106 g/mol (5,12$\cdot$53/(1,28$\cdot$2)); A = p{-}xileno (1,4{-}dimetilbenzeno), C8H10. Sua monohalogenação no anel gera um único produto, o 2{-}X{-}1,4{-}dimetilbenzeno. & M = 106 g/mol; A = p{-}xileno (1,4{-}dimetilbenzeno) & igual & ambos \\
\bottomrule\end{longtable}%
\begin{itemize}%
\item%
\textbf{unico)} Conferi o enunciado no PDF original (pág. 5): 53 g de A, 2 kg de benzeno, $\Delta$Tf = 1,28 $^\circ$C, halogenação com um único produto possível e Tollens negativo. A transcrição está fiel. Refiz a conta: M = Kf$\cdot$m/($\Delta$Tf$\cdot$m\_solv) = 5,12$\times$53/(1,28$\times$2) = 271,36/2,56 = 106,0 g/mol. Com massas atômicas precisas, C8H10 dá 106,17 e C7H6O dá 106,12, ambos compatíveis. A lógica é a clássica da questão: com M = 106 os candidatos são o benzaldeído, o etilbenzeno e os três xilenos. O Tollens negativo exclui o benzaldeído. Na substituição eletrofílica no anel, o etilbenzeno e o m{-}xileno dão 3 produtos e o o{-}xileno dá 2. Só o p{-}xileno, com os 4 H do anel equivalentes, dá um único produto monohalogenado. A resolução e o gabarito chegam à mesma massa molar e à mesma estrutura, e os dois estão certos.%
\end{itemize}%
\paragraph{Pontos de dúvida}%
\begin{itemize}%
\item%
\textit{(solution)} No Passo 2, a resolução afirma que 'um anel fundido ao benzeno exigiria pelo menos C9 (108 g/mol)'. Isso é impreciso: o benzociclobuteno (C8H8) tem só 8 carbonos e M = 104 g/mol. A conclusão não muda, porque 104 fica fora da janela de 105 a 107, mas a justificativa correta seria comparar com as massas reais dos bicíclicos, e não com um mínimo de C9.%
\item%
\textit{(solution)} A tabela de halogenação mistura critérios. Para quase todos os compostos, a coluna 'Classes de H no anel' conta posições não equivalentes por simetria (por exemplo, o, m e p no etilbenzeno). Para a N{-}metilanilina, porém, lista só 'o, p', que é o efeito dirigente e não a simetria. Pelo critério de simetria, seriam três classes. Isso não afeta a eliminação, porque a N{-}metilanilina dá mais de um produto de qualquer jeito.%
\item%
\textit{(statement)} O enunciado não diz se a halogenação é no anel (SEAr com FeX3) ou na cadeia lateral (radicalar), nem se é monossubstituição. A resolução declara a hipótese (H5) e a justifica: na via radicalar, os três xilenos dariam um único ArCH2X cada, e a questão não teria resposta única. Pela via do anel, a resposta p{-}xileno fica bem determinada. É uma hipótese razoável e coincide com a leitura do gabarito.%
\item%
\textit{(solution)} As hipóteses i = 1 (sem associação em benzeno) e 'aromático = derivado benzênico' (H4) são implícitas no enunciado. A tropona (C7H6O, 106) e heteroaromáticos como as lutidinas (107) ficam excluídos por H4. Mesmo sem H4, nenhum deles daria um único produto de halogenação, então a resposta não muda.%
\end{itemize}%
\section*{Questão 10}%
\href{https://poliedroresolve.sistemapoliedro.com.br/vestibulares/ime/2027/ime-2-fase-dia-3-28-10-2026-2-fase-dia-3/dissertativa-10-quimica-2-fase-dia-3-ime-2027}{Resolução do gabarito}\par\smallskip%
\begin{longtable}{@{}p{0.06\textwidth} p{0.33\textwidth} p{0.33\textwidth} p{0.1\textwidth} p{0.08\textwidth}@{}}\toprule Item & Resolução & Gabarito & Veredito & Certo \\ \midrule
unico & t = (T$\frac12$/ln2)$\cdot$ln(1+R). O valor exato dá $\approx$ 4,83$\times$10$^{8}$ anos; a linearização com os ln da Folha de Dados dá $\approx$ 4,88$\times$10$^{8}$ anos. Resposta declarada: 4,8 a 4,9$\times$10$^{8}$ anos (cerca de 5$\times$10$^{8}$ anos). & t/t$\frac12$ = 0,03463 $\Rightarrow$ t $\approx$ 4,88$\cdot$10$^{8}$ anos & equivalente & ambos \\
\bottomrule\end{longtable}%
\begin{itemize}%
\item%
\textbf{unico)} O modelo é o mesmo nos dois: N\_Pb/N\_Th = e\^{}\{$\lambda$t\} - 1 = 0,024, com $\lambda$ = ln2/T$\frac12$($^2$$^3$$^2$Th). Refazendo a conta: o gabarito usa t/t$\frac12$ = 0,024/0,693 = 0,034632. Isso equivale a ln(1,024) $\approx$ 0,024, que também sai da Folha de Dados como 7$\cdot$0,693 - 3$\cdot$1,609 = 0,024. Daí t = 0,034632 $\times$ 1,41$\times$10$^{1}$$^{0}$ = 4,883$\times$10$^{8}$ anos, exatamente o valor 'linearizado' que a resolução calcula explicitamente. Com ln(1,024) exato (0,023717), t/t$\frac12$ = 0,03422 e t = 4,82--4,83$\times$10$^{8}$ anos (4,824 com ln2 exato; 4,826 com 0,693). A diferença é de cerca de 1,2\% ($\approx$ R/2), menor que a incerteza da razão 0,024: R entre 0,0235 e 0,0245 dá t entre 4,73 e 4,92$\times$10$^{8}$ anos. A faixa final da resolução (4,8--4,9$\times$10$^{8}$) contém o valor do gabarito, que aparece como uma das duas estimativas. As verificações auxiliares também estão corretas: a cadeia tem 6 $\alpha$ e 4 $\beta$$^{-}$; $\lambda$ $\approx$ 4,9$\times$10$^{-}$$^{1}$$^{1}$ ano$^{-}$$^{1}$; a soma dos intermediários é $\approx$ 5$\times$10$^{-}$$^{1}$$^{0}$ N\_Th; as armadilhas dão -7,6$\times$10$^{1}$$^{0}$, 7,6$\times$10$^{1}$$^{0}$, 1,5$\times$10$^{8}$ e 3,3$\times$10$^{8}$ anos.%
\end{itemize}%
\paragraph{Pontos de dúvida}%
\begin{itemize}%
\item%
\textit{(key)} O valor t/t$\frac12$ = 0,03463 do gabarito vem da aproximação ln(1,024) $\approx$ 0,024 (linearização, ou 7ln2 - 3ln5 com os valores da Folha de Dados). O valor exato é t/t$\frac12$ = 0,03422, que dá t $\approx$ 4,82$\times$10$^{8}$ anos. A diferença é de \textasciitilde{}1,2\% e está dentro da precisão dos dados, mas um corretor rígido poderia esperar 4,88$\times$10$^{8}$ e não 4,83$\times$10$^{8}$.%
\item%
\textit{(solution)} A resposta final vem como faixa (4,8 a 4,9$\times$10$^{8}$ anos) e não como valor único. A faixa é justificada pela incerteza de R e contém explicitamente 4,88$\times$10$^{8}$, então não deve tirar ponto. Ainda assim, vale confirmar se a banca aceita resposta em faixa. O valor 'exato' 4,83$\times$10$^{8}$ usa ln2 $\approx$ 0,693; com ln2 exato fica 4,82$\times$10$^{8}$, diferença irrelevante.%
\item%
\textit{(statement)} O enunciado dá dados que não são usados (9\% de ThO$_{2}$, 0,35\% de U$_{3}$O$_{8}$, $^2$$^{0}$$^{6}$Pb e T$\frac12$ do $^2$$^3$$^{8}$U) e não dá a razão $^2$$^{0}$$^{6}$Pb/$^2$$^3$$^{8}$U. Só a série do $^2$$^3$$^2$Th fornece a idade. A resolução trata isso corretamente.%
\end{itemize}%
\end{document}